\section{Method}
\subsection{Problem Formulation}
Let $\mathcal{L}=\{(X_i,Y_i)\}_{i=1}^{N}$ denote the labeled set and $\mathcal{U}=\{U_i\}_{i=1}^{M}$ denote the unlabeled set, where $M \geq N$. Each image is $X_i,U_i\in\mathbb{R}^{W\times H\times L}$, and each mask is $Y_i\in\{0,1,\ldots,C\}^{W\times H}$, where class 0 is background. Training data are sampled from $K$ domains. In MiDSS, $\mathcal{L}$ comes from one domain, while $\mathcal{U}$ contains samples from multiple domains without requiring domain labels \cite{ma2024constructing}.

The base semi-supervised training follows a Mean Teacher-style design \cite{tarvainen2017mean}. For an unlabeled image $U$, weak and strong augmentations generate $U_w$ and $U_s$. A teacher model predicts a probability map $\hat{P}_w$ on $U_w$, and the pseudo-label is $\hat{Q}_w=\arg\max(\hat{P}_w)$ \cite{lee2013pseudo}. To connect labeled and unlabeled domains, an intermediate sample is generated by copy-paste:
\begin{equation}
\begin{gathered}
U_c=X_w\odot M+U_s\odot(\mathbf{1}-M),\\
\hat{Q}_c=Y_w\odot M+\hat{Q}_w\odot(\mathbf{1}-M),
\end{gathered}
\label{eq:copy_paste}
\end{equation}
where $M\in\{0,1\}^{W\times H}$ is the pasted region mask and $\odot$ denotes element-wise multiplication.

\subsection{Overall Pipeline}
Fig.~\ref{fig:framework} shows the reproduced synergistic training pipeline. The two student models are U-Net, denoted $f^{UT}$, and MedSAM, denoted $f^{MS}$. Each has a corresponding teacher model, $\hat{f}^{UT}$ and $\hat{f}^{MS}$, updated by exponential moving average. For MedSAM, LoRA is applied to reduce the adaptation cost while preserving pretrained knowledge \cite{hu2021lora}.

\begin{figure*}[t]
\centering
\includegraphics[width=0.86\textwidth]{Framework.pdf}
\caption{Reproduced SynFoC pipeline. U-Net and MedSAM both produce pseudo-labels, and the two branches are trained jointly so that each can correct the other at different stages.}
\label{fig:framework}
\end{figure*}

The supervised losses of U-Net and MedSAM are computed independently on labeled images. For unlabeled images, the teacher predictions are first combined by an instance-wise ensemble ratio $\alpha$ to produce an ensembled probability map:
\begin{equation}
\hat{P}^{EN}_w=\alpha\hat{P}^{UT}_w+(1-\alpha)\hat{P}^{MS}_w,
\label{eq:ensemble}
\end{equation}
where $\alpha$ controls the contribution from U-Net. The ensembled pseudo-label is then inserted into the copy-paste sample $U_c$ to supervise both student branches, while consensus-divergence consistency regularization is applied to the predictions on $U_c$. The total objective is
\begin{equation}
L_{total}=L_x+\lambda(t)(L_u+L_c+L_d),
\label{eq:total_loss}
\end{equation}
where $L_x$ is the supervised loss, $L_u$ is the unsupervised pseudo-label loss, $L_c$ and $L_d$ are the consensus and divergence regularizers, and
\begin{equation}
\lambda(t)=e^{-5(1-t/t_{max})^2}
\label{eq:lambda}
\end{equation}
is the Gaussian ramp-up weight used during training.

\subsection{Self-Mutual Confidence}
The ensemble ratio should reflect whether U-Net has become reliable for a particular unlabeled instance. Let $Q^{UT}_w=\arg\max(P^{UT}_w)$, $\hat{Q}^{UT}_w=\arg\max(\hat{P}^{UT}_w)$, and $\hat{Q}^{MS}_w=\arg\max(\hat{P}^{MS}_w)$. SMC measures the reliability of the U-Net prediction from two perspectives.

\textbf{Self-confidence.} Self-confidence $\Phi^{self}$ measures the consistency between the U-Net student and teacher predictions:
\begin{equation}
\Phi^{self}=\frac{1}{C}\sum_{i=1}^{C}\frac{2\left|\mathds{1}(Q^{UT}_w=i)\cap\mathds{1}(\hat{Q}^{UT}_w=i)\right|}{\left|\mathds{1}(Q^{UT}_w=i)\right|+\left|\mathds{1}(\hat{Q}^{UT}_w=i)\right|},
\label{eq:self_conf}
\end{equation}
where $\mathds{1}(\cdot)$ is the indicator function and $|\cdot|$ counts the number of pixels in the corresponding binary mask.

\textbf{Mutual confidence.} Mutual confidence $\Phi^{mut}$ measures the agreement between the U-Net teacher and the MedSAM teacher:
\begin{equation}
\Phi^{mut}=\frac{1}{C}\sum_{i=1}^{C}\frac{2\left|\mathds{1}(\hat{Q}^{MS}_w=i)\cap\mathds{1}(\hat{Q}^{UT}_w=i)\right|}{\left|\mathds{1}(\hat{Q}^{MS}_w=i)\right|+\left|\mathds{1}(\hat{Q}^{UT}_w=i)\right|}.
\label{eq:mut_conf}
\end{equation}
When both values are high, the U-Net pseudo-label is treated as reliable and should receive a larger ensemble weight. Therefore, the instance-wise ratio is defined as
\begin{equation}
\alpha=\Phi^{self}\times\Phi^{mut}.
\label{eq:alpha}
\end{equation}

After computing $\alpha$, we obtain the ensembled pseudo-label $\hat{Q}^{EN}_w=\arg\max(\hat{P}^{EN}_w)$ and form the mixed pseudo-label
\begin{equation}
\hat{Q}^{EN}_c=Y_w\odot M+\hat{Q}^{EN}_w\odot(\mathbf{1}-M).
\label{eq:ensemble_copy_paste}
\end{equation}
We further keep only confident unlabeled regions using the mask
\begin{equation}
\begin{gathered}
W^{EN}_w=\mathds{1}(\max(\hat{P}^{EN}_w)\geq \tau),\\
W^{EN}_c=M+W^{EN}_w\odot(\mathbf{1}-M),
\end{gathered}
\label{eq:confidence_mask}
\end{equation}
where $\tau=0.95$ is the confidence threshold. The unsupervised loss on the copy-paste sample is then
\begin{equation}
\begin{aligned}
L_u={}&L_{ce}(\hat{Q}^{EN}_c,P^{UT}_c,W^{EN}_c)+L_{dice}(\hat{Q}^{EN}_c,P^{UT}_c,W^{EN}_c)\\
&+L_{ce}(\hat{Q}^{EN}_c,P^{MS}_c,W^{EN}_c)+L_{dice}(\hat{Q}^{EN}_c,P^{MS}_c,W^{EN}_c),
\end{aligned}
\label{eq:unsup_loss}
\end{equation}
where $P^{UT}_c$ and $P^{MS}_c$ are the student predictions on $U_c$.

\subsection{Consensus-Divergence Consistency Regularization}
SMC improves pseudo-label quality, but the two models can still disagree in difficult regions. CDCR separates the student predictions on $U_c$ into consensus and divergence regions:
\begin{equation}
\begin{gathered}
M_c=\mathds{1}(Q^{UT}_c=Q^{MS}_c),\\
M_d=\mathbf{1}-M_c,
\end{gathered}
\label{eq:cdcr_masks}
\end{equation}
where $Q^{UT}_c=\arg\max(P^{UT}_c)$ and $Q^{MS}_c=\arg\max(P^{MS}_c)$. Fig.~\ref{fig:modules} illustrates the SMC and CDCR modules.

\begin{figure}[t]
\centering
\includegraphics[width=0.95\linewidth]{modules.pdf}
\caption{SMC estimates pseudo-label reliability, while CDCR applies different constraints to agreement and disagreement regions.}
\label{fig:modules}
\end{figure}

\textbf{Consensus regularization.} In consensus regions, the models are encouraged to make confident predictions by minimizing entropy:
\begin{equation}
L_c=-\frac{1}{2S}\sum\left[(P^{UT}_c\log P^{UT}_c)\odot M_c+(P^{MS}_c\log P^{MS}_c)\odot M_c\right],
\label{eq:consensus_loss}
\end{equation}
where $S=W\times H\times C$. In practice, interpolation is used wherever the resolutions of U-Net and MedSAM predictions need to be aligned.

\textbf{Divergence regularization.} In divergence regions, the models are encouraged to reduce prediction discrepancies through mean squared error:
\begin{equation}
L_d=\frac{1}{S}\sum \mathrm{MSE}(P^{UT}_c,P^{MS}_c)\odot M_d.
\label{eq:divergence_loss}
\end{equation}
This separation lets the method reward agreement where the two branches already align and focus discrepancy reduction on the regions where their predictions differ most.
